\documentclass{beamer}

\usepackage{amsmath,amssymb,amsfonts,bm}
\usepackage{physics}
\usepackage{braket}

\title{Entanglement-Sensitive Observables for a Two-Qubit Quantum Sensor}
\subtitle{Spin and Spatial Entanglement, Collective Observables, and Links to High-Energy Physics}
\author{}
\date{}

\begin{document}

%================================================
\begin{frame}
  \titlepage
\end{frame}

%================================================
\begin{frame}{Motivation}

Consider two qubits used as a prototype quantum sensor.
The qubits may be entangled in
\begin{itemize}
    \item spin,
    \item spatial mode,
    \item or both.
\end{itemize}

A natural question is:
\begin{quote}
Which observables best reveal the role of entanglement in the sensing protocol?
\end{quote}

A second question is whether one can go beyond local measurements and use
observables acting on the whole two-qubit system.

\end{frame}

%================================================
\begin{frame}{Basic Setup}

Let the two qubits be labeled $A$ and $B$.
For spin-$1/2$ degrees of freedom, define Pauli operators
\[
\sigma_x^{(A)},\sigma_y^{(A)},\sigma_z^{(A)},\qquad
\sigma_x^{(B)},\sigma_y^{(B)},\sigma_z^{(B)}.
\]

We may also define collective operators
\[
J_\alpha = \frac{1}{2}\left(\sigma_\alpha^{(A)}+\sigma_\alpha^{(B)}\right),
\qquad \alpha=x,y,z.
\]

If one qubit couples to a magnetic field $B$, a simple interaction Hamiltonian is
\[
H_{\mathrm{int}} = \frac{\gamma B}{2}\sigma_z^{(A)}.
\]

A more symmetric model is
\[
H_{\mathrm{int}} = \frac{\gamma_A B_A}{2}\sigma_z^{(A)}+
\frac{\gamma_B B_B}{2}\sigma_z^{(B)}.
\]

\end{frame}

%================================================
\begin{frame}{Canonical Entangled Sensor States}

Typical two-qubit entangled states include the Bell states
\begin{align}
\ket{\Phi^\pm} &= \frac{1}{\sqrt{2}}(\ket{00}\pm\ket{11}),\\
\ket{\Psi^\pm} &= \frac{1}{\sqrt{2}}(\ket{01}\pm\ket{10}).
\end{align}

For sensing, these states are important because they respond differently to
common-mode and differential fields.

Examples:
\begin{itemize}
    \item $\ket{\Phi^\pm}$ are sensitive to collective phase accumulation,
    \item $\ket{\Psi^\pm}$ can be sensitive to phase gradients and differential couplings.
\end{itemize}

\end{frame}

%================================================
\begin{frame}{Local Observables}

The most obvious observables are local spin components:
\[
\expval{\sigma_\alpha^{(A)}}, \qquad \expval{\sigma_\alpha^{(B)}}.
\]

If only qubit $A$ is exposed to the signal field, then
\[
\expval{\sigma_z^{(A)}}, \quad \expval{\sigma_x^{(A)}}, \quad \expval{\sigma_y^{(A)}}
\]
can reveal local phase accumulation.

However, these local observables do not always capture the metrological value
of entanglement. In many cases, entanglement is invisible in all one-body marginals
but appears clearly in correlations.

\end{frame}

%================================================
\begin{frame}{Two-Body Correlation Observables}

The simplest genuinely nonlocal observables are correlators
\[
C_{\alpha\beta}=\expval{\sigma_\alpha^{(A)}\sigma_\beta^{(B)}},
\qquad \alpha,\beta\in\{x,y,z\}.
\]

Important examples:
\begin{align}
\expval{\sigma_z^{(A)}\sigma_z^{(B)}}, \\
\expval{\sigma_x^{(A)}\sigma_x^{(B)}}, \\
\expval{\sigma_y^{(A)}\sigma_y^{(B)}}.
\end{align}

These observables:
\begin{itemize}
    \item distinguish Bell states,
    \item quantify alignment/anti-alignment,
    \item detect how phase information is encoded jointly in the two qubits.
\end{itemize}

They are often the most direct observables for testing whether entanglement improves sensing.

\end{frame}

%================================================
\begin{frame}{Connected Correlators}

A better diagnostic than the raw correlator is often the connected correlation
\[
\langle \sigma_\alpha^{(A)}\sigma_\beta^{(B)} \rangle_c
=
\expval{\sigma_\alpha^{(A)}\sigma_\beta^{(B)}}
-
\expval{\sigma_\alpha^{(A)}}\expval{\sigma_\beta^{(B)}}.
\]

Why useful?
\begin{itemize}
    \item It removes trivial product-state contributions.
    \item It isolates genuine correlation effects.
    \item It is often nonzero precisely when entanglement is operationally relevant.
\end{itemize}

In sensing, a strong field dependence of connected correlators can signal
that the parameter is encoded nonlocally.

\end{frame}

%================================================
\begin{frame}{Collective Spin Observables}

Observables acting on the whole system include
\[
J_x,\qquad J_y,\qquad J_z
\]
and quadratic forms
\[
J_x^2,\qquad J_y^2,\qquad J_z^2.
\]

For two qubits,
\[
J_\alpha^2 = \frac{1}{2}\left(I+\sigma_\alpha^{(A)}\sigma_\alpha^{(B)}\right).
\]

Thus collective spin fluctuations directly probe two-body correlations.

Important quantities:
\begin{align}
\expval{J_\alpha}, \\
(\Delta J_\alpha)^2 = \expval{J_\alpha^2}-\expval{J_\alpha}^2.
\end{align}

These are natural observables for testing whether entanglement sharpens or broadens phase sensitivity.

\end{frame}

%================================================
\begin{frame}{Parity Observables}

A very important global observable is parity.
For example, along the $z$ axis:
\[
\Pi_z = \sigma_z^{(A)}\sigma_z^{(B)}.
\]

More generally, after a collective rotation one can measure
\[
\Pi(\phi) = U^\dagger(\phi)\,\sigma_z^{(A)}\sigma_z^{(B)}\,U(\phi).
\]

Parity is central in quantum metrology because
\begin{itemize}
    \item it is highly sensitive to phase accumulation,
    \item it can saturate optimal sensitivity in some entangled interferometers,
    \item it is directly measurable in many qubit platforms.
\end{itemize}

For Bell-like states, parity oscillations are often a clean signature of entanglement-enhanced sensing.

\end{frame}

%================================================
\begin{frame}{Bell Operators and CHSH-Type Observables}

Another nonlocal observable class is the Bell operator, e.g.
\[
\mathcal{B}
=
\mathbf{a}\cdot\bm{\sigma}^{(A)}\otimes
(\mathbf{b}+\mathbf{b}')\cdot\bm{\sigma}^{(B)}
+
\mathbf{a}'\cdot\bm{\sigma}^{(A)}\otimes
(\mathbf{b}-\mathbf{b}')\cdot\bm{\sigma}^{(B)}.
\]

Its expectation value probes Bell nonlocality.

Why relevant for sensing?
\begin{itemize}
    \item Bell-operator expectation values are built from precisely the correlators
    that encode nonclassical joint response.
    \item If the signal changes $\expval{\mathcal{B}}$, then the sensing process is affecting
    nonlocal quantum correlations, not merely local polarization.
\end{itemize}

This is especially useful when studying robustness of sensing advantage against decoherence.

\end{frame}

%================================================
\begin{frame}{Quantum Fisher Information as the Central Quantity}

For parameter estimation, the most important quantity is often the
quantum Fisher information (QFI),
\[
F_Q[\rho,G],
\]
where $G$ is the generator of the parameter encoding.

For pure states,
\[
F_Q = 4(\Delta G)^2.
\]

Thus the observable content is encoded in the variance of the generator.

Examples:
\begin{itemize}
    \item If the field couples only to qubit $A$, then $G=\sigma_z^{(A)}/2$.
    \item If the field couples collectively, then $G=J_z$.
    \item If one probes a gradient field, then $G=(\sigma_z^{(A)}-\sigma_z^{(B)})/2$.
\end{itemize}

Entanglement helps when it enhances the variance of the relevant global generator.

\end{frame}

%================================================
\begin{frame}{Observables for Uniform and Gradient Fields}

Two especially useful whole-system generators are

\[
G_{\mathrm{com}} = \frac{1}{2}\left(\sigma_z^{(A)}+\sigma_z^{(B)}\right)=J_z
\]
for common-mode fields, and
\[
G_{\mathrm{diff}} = \frac{1}{2}\left(\sigma_z^{(A)}-\sigma_z^{(B)}\right)
\]
for gradient or differential fields.

Then one studies
\[
(\Delta G_{\mathrm{com}})^2,\qquad
(\Delta G_{\mathrm{diff}})^2.
\]

Different entangled states optimize different sensing tasks:
\begin{itemize}
    \item $\ket{\Phi^\pm}$ can be highly sensitive to common phase shifts,
    \item $\ket{\Psi^\pm}$ can be highly sensitive to differential phase shifts.
\end{itemize}

\end{frame}

%================================================
\begin{frame}{If the Qubits Are Also Spatially Entangled}

Suppose each qubit also has a spatial degree of freedom, for example
two paths or two locations.
Then we may define mode operators or effective spatial Pauli operators
\[
\tau_x,\tau_y,\tau_z
\]
acting on the spatial subspace.

Then one can study hybrid observables such as
\[
\sigma_z^{(A)}\otimes \tau_z^{(A)},\qquad
\sigma_z^{(A)}\sigma_z^{(B)},\qquad
\tau_x^{(A)}\tau_x^{(B)},
\]
or mixed spin-space correlators
\[
\sigma_x^{(A)}\tau_x^{(A)}\,\sigma_x^{(B)}\tau_x^{(B)}.
\]

These observables probe whether the sensing advantage arises from
\begin{itemize}
    \item spin entanglement,
    \item spatial entanglement,
    \item or spin-space hybrid entanglement.
\end{itemize}

\end{frame}

%================================================
\begin{frame}{Total-Spin Observables}

For two spin-$1/2$ qubits, another very natural observable is
\[
\mathbf{J}^2 = J_x^2 + J_y^2 + J_z^2.
\]

This distinguishes singlet and triplet sectors:
\begin{align}
\mathbf{J}^2 \ket{\text{triplet}} &= 2 \ket{\text{triplet}},\\
\mathbf{J}^2 \ket{\text{singlet}} &= 0.
\end{align}

This is useful because:
\begin{itemize}
    \item the singlet state is rotationally invariant,
    \item triplet states respond differently to external fields,
    \item total-spin measurements reveal how symmetry constrains sensing response.
\end{itemize}

This observable acts on the whole two-qubit system and is often more informative
than local spin projections alone.

\end{frame}

%================================================
\begin{frame}{Entanglement Witnesses as Observables}

One may also construct entanglement witnesses $W$ such that
\[
\expval{W} < 0
\]
signals entanglement.

For example, witnesses built from correlators
\[
W \sim I - \sigma_x^{(A)}\sigma_x^{(B)}
- \sigma_y^{(A)}\sigma_y^{(B)}
- \sigma_z^{(A)}\sigma_z^{(B)}
\]
can be directly measurable.

These are not always optimal sensing observables, but they are useful for
\begin{itemize}
    \item verifying that the enhanced sensitivity truly comes from entanglement,
    \item separating entanglement-enhanced performance from classical correlation effects.
\end{itemize}

\end{frame}

%================================================
\begin{frame}{Which Observables Are Best for Testing Entanglement in Sensing?}

A practical hierarchy is:

\begin{enumerate}
    \item Local magnetization:
    \[
    \expval{\sigma_z^{(A)}},\quad \expval{\sigma_z^{(B)}}
    \]
    \item Two-body correlators:
    \[
    \expval{\sigma_\alpha^{(A)}\sigma_\beta^{(B)}}
    \]
    \item Collective observables:
    \[
    \expval{J_\alpha},\quad (\Delta J_\alpha)^2,\quad \expval{\mathbf{J}^2}
    \]
    \item Parity:
    \[
    \expval{\sigma_z^{(A)}\sigma_z^{(B)}}
    \]
    \item Generator variances / QFI:
    \[
    (\Delta G)^2,\quad F_Q
    \]
\end{enumerate}

If the goal is specifically to quantify metrological gain, then
generator variance and QFI are the most directly relevant quantities.

\end{frame}

%================================================
\begin{frame}{Relevance for High-Energy Physics}

Yes, several of these observables are relevant in high-energy physics.

Examples include:
\begin{itemize}
    \item spin-correlation observables in relativistic two-particle production,
    \item Bell-type correlators for particle-antiparticle systems,
    \item flavor and spin oscillation observables,
    \item collective symmetry generators in quantum field theory,
    \item detector-response observables in relativistic quantum information.
\end{itemize}

The same mathematical structures appear:
\[
\expval{\sigma_\alpha^{(A)}\sigma_\beta^{(B)}},\qquad
\expval{\mathcal{B}},\qquad
(\Delta G)^2.
\]

Thus the language of quantum sensing overlaps strongly with the language of
correlation measurements in high-energy and relativistic physics.

\end{frame}

%================================================
\begin{frame}{Examples Relevant to High-Energy Physics}

Some concrete directions where these observables matter:

\begin{itemize}
    \item \textbf{Top-quark or lepton spin correlations:}
    measured through correlation tensors analogous to
    \[
    \expval{\sigma_i^{(A)}\sigma_j^{(B)}}.
    \]

    \item \textbf{Neutral meson or particle-antiparticle entanglement:}
    Bell-like observables and parity-type observables test quantum coherence.

    \item \textbf{Neutrino oscillation and flavor interferometry:}
    collective generators and phase observables resemble sensing generators.

    \item \textbf{Unruh/Hawking detector models:}
    two-detector correlations and entanglement harvesting use very similar
    joint observables.
\end{itemize}

So the same observables that diagnose entanglement-enhanced sensing can also be
reinterpreted as probes of relativistic quantum correlations.

\end{frame}

%================================================
\begin{frame}{A Useful Minimal Measurement Program}

For a two-qubit entangled quantum sensor, a minimal but powerful set of measurements is:

\begin{align}
\expval{\sigma_z^{(A)}}, \qquad
\expval{\sigma_z^{(B)}}, \\
\expval{\sigma_x^{(A)}\sigma_x^{(B)}}, \qquad
\expval{\sigma_y^{(A)}\sigma_y^{(B)}}, \qquad
\expval{\sigma_z^{(A)}\sigma_z^{(B)}}, \\
\expval{J_z}, \qquad
(\Delta J_z)^2, \qquad
(\Delta G_{\mathrm{diff}})^2.
\end{align}

If spatial entanglement is also present, add
\[
\expval{\tau_x^{(A)}\tau_x^{(B)}},\qquad
\expval{\tau_z^{(A)}\tau_z^{(B)}},
\]
and mixed spin-space correlators.

This set already gives a strong picture of whether the advantage comes from
local response, joint response, or hybrid spin-space entanglement.

\end{frame}

%================================================
\begin{frame}{Summary}

For two spin and spatially entangled qubits used as a quantum sensor:

\begin{itemize}
    \item Local observables probe direct coupling to the field.
    \item Two-body correlators reveal genuinely entangled response.
    \item Collective observables such as $J_\alpha$, $J_\alpha^2$, and $\mathbf{J}^2$
    act on the whole system and are highly informative.
    \item Parity and Bell-type observables are powerful nonlocal diagnostics.
    \item The most fundamental metrological quantity is the QFI,
    determined by variances of the encoding generator.
\end{itemize}

Many of these observables are also directly relevant in high-energy physics,
especially in studies of spin correlations, Bell tests, relativistic detector models,
and quantum-field-theoretic interferometry.

\end{frame}

%================================================
\begin{frame}{Take-Home Message}

If entanglement is central to the sensing protocol, then measuring only
\[
\expval{\sigma_z^{(A)}}
\]
is usually not enough.

The most informative whole-system observables are often
\[
\expval{\sigma_\alpha^{(A)}\sigma_\beta^{(B)}},\qquad
\expval{J_\alpha},\qquad
\expval{\mathbf{J}^2},\qquad
\expval{\Pi},\qquad
F_Q.
\]

These capture how the parameter is encoded collectively and whether
entanglement genuinely improves sensitivity.

\end{frame}

\end{document}
